% GENERATED FILE. Edit canonical lessons, then run npm run generate:books.
% canonical-source-hash: fnv1a64:9f7344c0d3f285d7
% canonical-lessons: GE-C266-ploetzlich, GE-C266-unerwartet, GE-C266-allmaehlich, GE-C266-rechtzeitig, GE-C266-zugleich

\chapter{Auf einmal}
\label{ch:ge-auf-einmal}
\hlchaptermodality{266}

\begin{chapteropening}
\textbf{By the end of this chapter:} \emph{I can say how fast things happened: suddenly, unexpectedly, gradually, in time and at once.}

Complete the last lesson of chapter 266: I can say how fast things happened: suddenly, unexpectedly, gradually, in time and at once.
\end{chapteropening}

\section[plötzlich]{plötzlich — suddenly}

\label{lesson:GE-C266-ploetzlich}



\begin{quote}
Before the new one: say the German for so far, then the German for long ago.
\end{quote}

\subsection*{You'll want to know: plötzlich}
\textbf{plötzlich} — \textquotedblleft{}suddenly\textquotedblright{}.

\begin{quote}
Plötzlich ging das Licht aus.
\end{quote}

\begin{quote}
\emph{Suddenly the light went out.}
\end{quote}

Stories love it; \emph{auf einmal} means the same. Also an adjective: \emph{ein plötzlicher Regen}.

\begin{cousinweb}
\emph{plötzlich} is usually traced to Middle High German \emph{plotz} \textquotedblleft{}a clap, a slap\textquotedblright{} (EtymDB), a word that imitates the sound of a thud; EtymDB lists no parent for \emph{plötzlich} itself.
\end{cousinweb}

\begin{grammarlens}[title={What you've built}]
One more word for telling, explaining and describing at B1.
\end{grammarlens}

\subsection*{Your turn}
\begin{itemize}
\raggedright
  \item \emph{Say it:} \emph{plötzlich}
  \item \emph{Say it:} \emph{plötzlich}, once more
  \item \emph{Say it:} a sentence of your own with \emph{plötzlich}
  \item \emph{Recall:} say the German for so far, then the German for long ago, then say \emph{plötzlich} again
\end{itemize}

\begin{tcolorbox}[breakable,colback=teal!4,colframe=teal!35!black,title={Before you move on}]
What does plötzlich mean? (\textquotedblleft{}Suddenly\textquotedblright{}.) Say it once more.
\end{tcolorbox}
\section[unerwartet]{unerwartet — unexpectedly; unexpected}

\label{lesson:GE-C266-unerwartet}



\begin{quote}
Before the new one: say the German for long ago, then the German for suddenly.
\end{quote}

\subsection*{You'll want to know: unerwartet}
\textbf{unerwartet} — \textquotedblleft{}unexpectedly; unexpected\textquotedblright{}.

\begin{quote}
Mein Bruder kam unerwartet zu Besuch.
\end{quote}

\begin{quote}
\emph{My brother came to visit unexpectedly.}
\end{quote}

The opposite of \emph{erwartet}, from \emph{erwarten} 'to expect'.

\begin{cousinweb}
\emph{unerwartet} joins \emph{un-} \textquotedblleft{}not\textquotedblright{} and \emph{erwartet} \textquotedblleft{}expected\textquotedblright{} (EtymDB). The base is \emph{warten}, from Old High German \emph{wartēn} \textquotedblleft{}to watch, to look after\textquotedblright{} (EtymDB).
\end{cousinweb}

\begin{grammarlens}[title={What you've built}]
One more word for telling, explaining and describing at B1.
\end{grammarlens}

\subsection*{Your turn}
\begin{itemize}
\raggedright
  \item \emph{Say it:} \emph{unerwartet}
  \item \emph{Say it:} \emph{unerwartet}, once more
  \item \emph{Say it:} a sentence of your own with \emph{unerwartet}
  \item \emph{Recall:} say the German for long ago, then the German for suddenly, then say \emph{unerwartet} again
\end{itemize}

\begin{tcolorbox}[breakable,colback=teal!4,colframe=teal!35!black,title={Before you move on}]
What does unerwartet mean? (\textquotedblleft{}Unexpectedly; unexpected\textquotedblright{}.) Say it once more.
\end{tcolorbox}
\section[allmählich]{allmählich — gradually}

\label{lesson:GE-C266-allmaehlich}



\begin{quote}
Before the new one: say the German for suddenly, then the German for unexpectedly.
\end{quote}

\subsection*{You'll want to know: allmählich}
\textbf{allmählich} — \textquotedblleft{}gradually\textquotedblright{}.

\begin{quote}
Allmählich wurde es dunkel.
\end{quote}

\begin{quote}
\emph{Gradually it got dark.}
\end{quote}

Close to \emph{nach und nach}.

\begin{cousinweb}
\emph{allmählich} comes from Middle High German \emph{almehlich} (EtymDB), usually traced to \emph{gemach} \textquotedblleft{}slow, easy\textquotedblright{}: little by little.
\end{cousinweb}

\begin{grammarlens}[title={What you've built}]
One more word for telling, explaining and describing at B1.
\end{grammarlens}

\subsection*{Your turn}
\begin{itemize}
\raggedright
  \item \emph{Say it:} \emph{allmählich}
  \item \emph{Say it:} \emph{allmählich}, once more
  \item \emph{Say it:} a sentence of your own with \emph{allmählich}
  \item \emph{Recall:} say the German for suddenly, then the German for unexpectedly, then say \emph{allmählich} again
\end{itemize}

\begin{tcolorbox}[breakable,colback=teal!4,colframe=teal!35!black,title={Before you move on}]
What does allmählich mean? (\textquotedblleft{}Gradually\textquotedblright{}.) Say it once more.
\end{tcolorbox}
\section[rechtzeitig]{rechtzeitig — in time, in good time}

\label{lesson:GE-C266-rechtzeitig}



\begin{quote}
Before the new one: say the German for unexpectedly, then the German for gradually.
\end{quote}

\subsection*{You'll want to know: rechtzeitig}
\textbf{rechtzeitig} — \textquotedblleft{}in time, in good time\textquotedblright{}.

\begin{quote}
Wir sind rechtzeitig am Bahnhof angekommen.
\end{quote}

\begin{quote}
\emph{We arrived at the station in time.}
\end{quote}

In time for something; \emph{pünktlich} means exactly at the time.

\begin{cousinweb}
\emph{rechtzeitig} joins \emph{recht} \textquotedblleft{}right\textquotedblright{} and \emph{Zeit} \textquotedblleft{}time\textquotedblright{}, both inherited Germanic words (EtymDB), with the ending \emph{-ig}: at the right time.
\end{cousinweb}

\begin{grammarlens}[title={What you've built}]
One more word for telling, explaining and describing at B1.
\end{grammarlens}

\subsection*{Your turn}
\begin{itemize}
\raggedright
  \item \emph{Say it:} \emph{rechtzeitig}
  \item \emph{Say it:} \emph{rechtzeitig}, once more
  \item \emph{Say it:} a sentence of your own with \emph{rechtzeitig}
  \item \emph{Recall:} say the German for unexpectedly, then the German for gradually, then say \emph{rechtzeitig} again
\end{itemize}

\begin{tcolorbox}[breakable,colback=teal!4,colframe=teal!35!black,title={Before you move on}]
What does rechtzeitig mean? (\textquotedblleft{}In time, in good time\textquotedblright{}.) Say it once more.
\end{tcolorbox}
\section[zugleich]{zugleich — at once, at the same moment; both}

\label{lesson:GE-C266-zugleich}



\begin{quote}
Before the new one: say the German for gradually, then the German for in time.
\end{quote}

\subsection*{You'll want to know: zugleich}
\textbf{zugleich} — \textquotedblleft{}at once, at the same moment; both\textquotedblright{}.

\begin{quote}
Alle sprachen zugleich.
\end{quote}

\begin{quote}
\emph{Everyone spoke at once.}
\end{quote}

Close to \emph{gleichzeitig}, and more written. Also 'both': \emph{Sie ist Ärztin und zugleich Mutter}.

\begin{cousinweb}
\emph{zugleich} joins \emph{zu} and \emph{gleich} \textquotedblleft{}equal, alike\textquotedblright{}, from Old High German \emph{gilīh} (EtymDB), usually traced to Proto-Germanic \emph{galīkaz} \textquotedblleft{}alike, similar\textquotedblright{}: in the same way, at the same time.
\end{cousinweb}

\begin{grammarlens}[title={What you've built}]
One more word for telling, explaining and describing at B1.
\end{grammarlens}

\subsection*{Your turn}
\begin{itemize}
\raggedright
  \item \emph{Say it:} \emph{zugleich}
  \item \emph{Say it:} \emph{zugleich}, once more
  \item \emph{Say it:} a sentence of your own with \emph{zugleich}
  \item \emph{Recall:} say the German for gradually, then the German for in time, then say \emph{zugleich} again
\end{itemize}

\begin{tcolorbox}[breakable,colback=teal!4,colframe=teal!35!black,title={Before you move on}]
What does zugleich mean? (\textquotedblleft{}At once, at the same moment; both\textquotedblright{}.) Say it once more.
\end{tcolorbox}
